% Responsable: P1
\section{Conclusiones}
\label{sec:conclusiones}

\paragraph{Una CFG describe bien un fragmento del español, y solo ese fragmento.} Con seis no
terminales y un léxico de 221 palabras, la gramática cubre oraciones con sujeto, verbo, objeto,
adjetivos, adverbios, frases preposicionales y dos tiempos verbales. Dentro de ese fragmento el
parser es exacto: acepta las 45 oraciones de la gramática, rechaza las 28 que no lo están por la
razón correcta, y etiqueta y encuentra el sujeto, el verbo y el objeto sin errores. Pero fuera de
él no da nada: el 20\,\% del español correcto de nuestro conjunto (sujeto omitido, orden libre,
coordinación, subordinadas, palabras nuevas) queda rechazado sin análisis.

\paragraph{La concordancia es donde la CFG con rasgos aporta lo que los demás no tienen.} Es el
único sistema que rechaza las 13 oraciones con error de concordancia, dice qué palabra falla y no
da falsas alarmas. Los etiquetadores las analizan como oraciones válidas, y la puntuación de un
modelo de lenguaje las rechaza junto con oraciones correctas poco frecuentes, sin poder decir
por qué. Separar la concordancia de la gramática, como la verificación semántica de la Fase~1, fue
lo que la hizo manejable: escribirla en las reglas habría multiplicado la gramática.

\paragraph{La ambigüedad es el límite común.} La gramática no es LL(1) por la frase
preposicional, y ninguna transformación lo arregla porque la ambigüedad es del idioma. Nuestro
parser la resuelve con una regla fija que acierta en el 59\,\% de los casos; spaCy y el LLM
aciertan más porque usan el sentido, pero ningún sistema reconoce las oraciones ambiguas como
tales.

\paragraph{Los modelos estadísticos ganan en cobertura y pierden en garantías.} XLM-R etiqueta
el 98\,\% de las palabras de todo el conjunto, y el 100\,\% de las que el parser puede analizar,
sin léxico ni reglas. Pero un etiquetador no decide la gramaticalidad, y la decisión de BETO
mezcla gramaticalidad con frecuencia.

\paragraph{Un LLM puede seguir reglas formales, pero no se puede confiar en que lo haga.} Con la
gramática en el prompt, el LLM decide casi como el parser (exactitud 0.97--0.99 contra las
reglas) en lugar de seguir su conocimiento del español. Pero justifica algunas decisiones con
árboles que no cumplen la gramática, y en el caso más claro inventa una derivación que no
corresponde a la oración. Los ejemplos del \emph{few-shot} mejoran mucho la forma de los árboles
y poco la decisión.

\paragraph{Limitaciones y trabajo futuro.}
\begin{itemize}
  \item El conjunto de prueba tiene 73 oraciones y se etiquetó con la misma gramática que
    implementa el parser; con más oraciones, y oraciones escritas por otras personas, las
    diferencias entre sistemas serían más confiables.
  \item El parser no retrocede: \emph{El gato bajo la mesa duerme} está en la gramática y se
    rechaza. Un parser con retroceso, o un parser de gráfico que devuelva todos los árboles,
    resolvería eso y permitiría mostrar las dos lecturas de las oraciones ambiguas.
  \item La gramática se puede ampliar con sujeto omitido ($S \to VP$, con el sujeto deducido del
    verbo) y coordinación, que son las construcciones fuera de la gramática más frecuentes.
  \item El experimento con el LLM usó un solo modelo pequeño; repetirlo con modelos más
    grandes, como \texttt{gemini-3.8-flash}, diría si los errores en los árboles son de este
    modelo o del enfoque.
  \item Un sistema mixto podría combinar lo mejor de cada enfoque: un etiquetador estadístico en
    lugar del léxico cerrado, y la gramática con rasgos para verificar la estructura y la
    concordancia sobre esas etiquetas.
\end{itemize}
